\begin{textsample}{POS Dim 2 – 2010\_c5 – Score -2.00 – 2010\_c5/enade\_2010\_enf\_d4ae7ada6c.txt}  \label{ex:f2_pos_331}
CAPÍTULO 5 DISTRIBUIÇÃO DOS CONCEITOS 5.1 PANORAMA NACIONAL DA DISTRIBUIÇÃO DOS CONCEITOS A Tabela 5.1 apresenta a quantidade e a distribuição de cursos de Enfermagem participantes do ENADE/2010 por faixa de conceito e Grande Região.
Observando -se os dados da Tabela 5.1 nota -se que, dos 691 cursos participantes, 206 ( 29,8 \% ) classificaram -se com conceito 2, o valor modal.
Este conceito modal, entretanto, foi o encontrado em três das cinco regiões brasileiras, com seguintes percentuais : 35,0 \% na região Norte, 34,4 \% na Sudeste e 27,1 \% na Centro-Oeste.
Nas regiões Nordeste ( 24,2 \% ) e Sul ( 34,6 \% ), o conceito modal foi o 3.
Houve, ainda, 28 ( 4,1 \% ) que receberam conceito 5 e 22 cursos ( 23,4 \% ) que receberam conceito 1, em âmbito nacional.
Dos 691 cursos, 162 ( 23,4 \% ) ficaram sem conceito ( SC ).
A região Norte participou com 40 cursos, 5,8 \% do total nacional.
Nesta região, o conceito 3 foi atribuído a sete cursos dos 40 participantes, equivalentes 17,5 \% do total regional.
Um pouco mais de ¼ dos cursos ( 27,5 \%, correspondendo a onze deles ) não receberam conceito.
Dois cursos receberam conceito 1, cinco o conceito 4 e somente um curso ( 2,5 \% ) recebeu conceito 5.
O conceito modal, como já mencionado, foi o 2, com 14 cursos.
A região Nordeste participou com 161 cursos, 23,3 \% do total nacional.
Destes, 39 cursos, 24,2 \% em termos regionais, obtiveram conceito 3.
Os conceitos 2 e 4 foram 5.2 CONCEITOS POR CATEGORIA ADMINISTRATIVA E POR GRANDE REGIÃO A Tabela 5.2 apresenta a distribuição dos cursos participantes do ENADE/2010 de Enfermagem, por Categoria Administrativa, de acordo com os conceitos por eles alcançados, segundo as Grandes Regiões brasileiras.
Dos 691 cursos participantes, 138 ( 20,0 \% ) eram ministrados em instituições públicas e 553 ( 80,0 \% ), em privadas.
De acordo com as informações da Tabela 5.2, as instituições públicas foram as que apresentaram, em termos nacionais, a maior proporção de cursos com conceito 5, equivalente a 15,9 \% ( 22 cursos ) do total dessa categoria.
Dos 138 cursos participantes em IES públicas, o conceito 4 foi atribuído a 53 cursos, o conceito modal.
Entre os demais cursos participantes, 22 apresentaram conceito 5, 20 o conceito 3, 11 o conceito 2 e seis apresentaram conceito 1.
Nesta categoria, 26 cursos ficaram sem conceito.
Na rede privada, o conceito modal foi 2 com 195 cursos.
Entre os demais cursos participantes, 176 receberam conceito 3 e 24, conceito 4.
O conceito 1 foi atribuído a 16 cursos e o conceito 5, a seis cursos.
Nesta categoria, 136 cursos ficaram sem conceito. atribuídos a, respectivamente, 30 e 22 cursos ( 18,6 \% e 13,7 \% ).
O conceito 1 foi atribuído a seis cursos ( 3,7 \% ), e o conceito 5, a dez cursos ( 6,2 \% ).
Nesta região, 33,5 \% dos cursos ( 54 cursos ) ficaram sem conceito.
Na região Sudeste, observou -se o conceito modal, 2 com 34,5 \%, correspondentes a 106 dos 307 cursos participantes na região.
O conceito 3 foi atribuído a 100 cursos ( 32,6 \% ) e o conceito 4, a 26 ( 8,5 \% ).
O conceito 1 foi a classificação obtida por sete cursos ( 2,3 \% ), e o conceito 5, a 14(4,6 \% ).
Do total de cursos da região, 54 ( 17,6 \% ) ficaram sem conceito.
Dos 113 cursos participantes da região Sul, 39 ( 34,5 \% ) obtiveram conceito 3, o conceito modal.
O conceito 2 foi atribuído a 37 cursos ( 32,7 \% ) e o conceito 4, a 15 ( 13,3 \% ).
Nenhum curso recebeu o conceito 5, ao passo que quatro obtiveram o conceito 1.
Os 18 cursos restantes ficaram sem conceito.
Um pouco mais que ¼ ( 27,1 \%, correspondendo a 19 cursos ) dos 70 cursos participantes na região Centro-Oeste recebeu conceito 2, o modal.
Receberam conceito 1 três cursos ( 4,3 \% ) e conceito 3, 11 cursos ( 15,7 \% ).
Nove cursos ( 12,9 \% ) foram avaliados com o conceito 4 e três ( 4,3 \% ), com o conceito 5.
Um pouco mais de 1/3 ( 35,7 \%, correspondendo a 25 cursos ) ficou sem conceito nesta região. 5.2 CONCEITOS POR CATEGORIA ADMINISTRATIVA E POR GRANDE REGIÃO A Tabela 5.2 apresenta a distribuição dos cursos participantes do ENADE/2010 de Enfermagem, por Categoria Administrativa, de acordo com os conceitos por eles alcançados, segundo as Grandes Regiões brasileiras.
Dos 691 cursos participantes, 138 ( 20,0 \% ) eram ministrados em instituições públicas e 553 ( 80,0 \% ), em privadas.
De acordo com as informações da Tabela 5.2, as instituições públicas foram as que apresentaram, em termos nacionais, a maior proporção de cursos com conceito 5, equivalente a 15,9 \% ( 22 cursos ) do total dessa categoria.
Dos 138 cursos participantes em IES públicas, o conceito 4 foi atribuído a 53 cursos, o conceito modal.
Entre os demais cursos participantes, 22 apresentaram conceito 5, 20 o conceito 3, 11 o conceito 2 e seis apresentaram conceito 1.
Nesta categoria, 26 cursos ficaram sem conceito.
Na rede privada, o conceito modal foi 2 com 195 cursos.
Entre os demais cursos participantes, 176 receberam conceito 3 e 24, conceito 4.
O conceito 1 foi atribuído a 16 cursos e o conceito 5, a seis cursos.
Nesta categoria, 136 cursos ficaram sem conceito.
Na análise por região, observa -se que, na região Norte, as instituições privadas participaram com 27 cursos, dos quais dez ficaram sem conceito.
O conceito modal para as instituições privadas na região foi 2, com 11 cursos.
Os demais foram avaliados com conceito 3 ( quatro cursos ) e conceito 1 ( dois cursos ).
As instituições públicas participaram com 13 cursos, dos quais somente um ficou sem conceito.
O conceito modal para esta Categoria Administrativa na região Norte foi 4 ( com cinco cursos ).
Os demais foram avaliados com conceito 2 ( três cursos ), conceito 3 ( três cursos ) e conceito 5 ( um curso ).
Na região Nordeste, a rede privada concentrou 110 dos 161 cursos participantes, 68,3 \% do total da região.
Entre os que receberam conceito nesta Categoria Administrativa na região Nordeste, predominaram os cursos com conceito 5 ( 44 cursos ).
Os demais foram avaliados com conceitos 1 ( 30 cursos ), 2 ( também 30 ) e 3 ( três cursos ), ou ficaram sem conceito ( três cursos ).
As instituições públicas nesta região participaram com 51 cursos, dos quais três ficaram sem conceito.
O conceito modal para esta Categoria Administrativa foi 3 ( com 19 cursos ).
Os demais foram avaliados com conceito 2 ( nove cursos ), conceito 4 ( dez cursos ) e 5 ( também dez ).
Na região Sudeste, a proporção de cursos da rede privada, 90,6 \%, foi mais elevada do que nas demais regiões brasileiras, correspondendo a 278 dos 307 cursos participantes.
Nesta categoria, na região Sudeste, o conceito modal foi 2 ( 102 cursos ).
Os demais foram avaliados com conceito 1 ( cinco cursos ), conceito 3 ( 98 cursos ), conceito 4 ( 18 cursos ) e conceito 5 ( seis cursos ).
Nesta combinação de Categoria Administrativa e Grande Região, 49 cursos não receberam conceito.
Entre os 29 cursos ( 9,4 \% ) em instituições públicas na região Sudeste, a categoria modal foi a que recebeu os conceitos 4 e 5 ( ambos com oito cursos ).
Os demais foram avaliados com conceito 1 ( dois cursos ), conceito 2 ( quatro cursos ), conceito 3 ( dois cursos ).
Nesta categoria na região Sudeste, cinco cursos ficaram sem conceito.
As instituições privadas concentraram 86 dos 113 cursos participantes da região Sul, 76,1 \% do total regional.
Desses, 35 foram avaliados no conceito modal, 3.
Os demais foram avaliados com os conceitos 1 ( quatro cursos ) e 2 ( 34 cursos ).
Nenhum curso de IES privadas, na região Sul, obteve os conceitos 4 ou 5.
Nesta combinação de Categoria Administrativa e Grande Região, 13 cursos não receberam conceito.
As instituições públicas, nesta região, participaram com 27 cursos, dos quais 15 receberam o conceito modal, 4.
Os demais foram avaliados com os conceitos 2 ( três cursos ) e 3 ( quatro cursos ).
Nesta categoria, na região Sul, cinco cursos ficaram sem conceito.
Na região Centro-Oeste, 52 dos 70 cursos participantes eram de instituições privadas ( 74,3 \% em termos regionais ).
Destes, 18 concentraram -se no conceito 2, o conceito modal, e 20 ficaram sem conceito.
Os demais receberam conceito 1 ( dois cursos ), conceito 3 ( nove cursos ) e 4 ( três deles ).
Na rede pública, dos 18 cursos na região Centro-Oeste ( 25,7 \% ), seis foram classificados com o conceito 4, o conceito modal, e cinco ficaram sem conceito.
Houve apenas um curso em cada, nos conceitos 1 e 2.
Dois receberam conceito 3, e três cursos, o conceito 5. 5.3 CONCEITOS POR ORGANIZAÇÃO ACADÊMICA E POR GRANDE REGIÃO Na Tabela 5.3 encontra -se a distribuição dos conceitos atribuídos aos cursos participantes do ENADE/2010 na Área de Enfermagem, por Organização Acadêmica, segundo as Grandes Regiões brasileiras.
Dos 691 cursos participantes, 270 eram oferecidos em Universidades, 89 em Centros Universitários, e os demais 332, em Faculdades.
De acordo com os dados apresentados, a grande maioria dos cursos avaliados com conceito 5 eram vinculados a Universidades ( 27 dos 28 com este conceito ).
Esse tipo de Organização Acadêmica teve o conceito 3 como modal, com 75 cursos.
Além dos 36 cursos que ficaram sem conceito, os demais classificaram -se com os conceitos 1 ( oito cursos ), 2 ( 60 cursos ) e 4 ( 64 cursos ).
Entre os cursos em Centros Universitários, o conceito modal também foi 3 com 39 cursos, e oito ficaram sem conceito.
Dos outros cursos neste tipo de Organização Acadêmica, dois deles receberam o conceito 1, enquanto 37, o conceito 2.
Três cursos foram classificados com o conceito 4 e nenhum atingiu o conceito 5.
Nas Faculdades, dos 332 cursos, 118 ficaram sem conceito e 109 receberam o conceito 2, o modal para este tipo de Organização Acadêmica entre os cursos com conceito.
Dos demais cursos, doze receberam conceito 1, 82 o conceito 3, dez o conceito 4 e apenas um o 5.
Considerando -se separadamente as regiões brasileiras, verifica -se que na região Norte as Universidades concentraram 14 dos 40 cursos participantes.
Dos cursos em Universidades, dois ficaram sem conceito e os doze cursos restantes apresentaram os seguintes conceitos : cinco cursos ficaram com o conceito 4 ( o conceito modal ), quatro cursos ficaram com conceito 3, dois com conceito 2 e um curso com conceito 5.
Os Centros Universitários foram representados por sete cursos, nenhum sem conceito.
Estes cursos receberam os conceitos 2 ( seis cursos, o conceito modal ) e 3 ( um curso ).
As Faculdades participaram com 19 cursos na região Norte.
Destes, nove ficaram sem conceito e os demais foram classificados da seguinte forma : conceito 1 ( dois cursos ), conceito 2 ( seis cursos ) e conceito 3 ( dois cursos ).
Na região Nordeste, as Universidades participaram com 57 cursos dos 161 na Área de Enfermagem da região.
Entre os cursos em Universidades, o conceito modal foi 4 ( 20 cursos ).
Os conceitos 3 e 5 foram a classificação de dez cursos para cada um dos conceitos.
Dos demais cursos, três cursos apresentaram conceito 1, ao passo que dois obtiveram conceito 2, enquanto 12 deles ficaram sem conceito.
Os Centros Universitários contaram com cinco cursos participantes na região.
O conceito modal foi 3 ( com três cursos ).
Os outros dois cursos receberam o conceito 2 ( dois cursos ).
As Faculdades foram representadas por 99 cursos, dos quais 42 não receberam conceito.
O conceito modal ficou entre o 2 e o 3, com 26 cursos em cada.
Dos cursos restantes, três receberam conceito 2 e dois receberam o 4.
Na região Sudeste, as Universidades concentraram 107 dos 307 cursos da região.
Entre os cursos em Universidades, o conceito modal foi 3, com 32 cursos, e 12 ficaram sem conceito.
Dos demais cursos, 31 receberam conceito 2, 16 o conceito 4, 13 o conceito 5 e três cursos receberam conceito 1.
Os Centros Universitários participaram com 59 cursos na região, dos quais 27 obtiveram o conceito modal, 3, e cinco ficaram sem conceito.
Os demais receberam conceito 2 ( 23 cursos ), conceito 4 ( três cursos ) e conceito 1 ( apenas um deles ).
As Faculdades foram representadas por 141 cursos na região, que se distribuíram nos conceitos 2 ( 52 cursos, o conceito modal ), conceito 3 ( 41 cursos ), conceito 4 ( sete cursos ), conceito 1 ( três cursos ) e conceito 5 ( apenas um deles ).
Ficaram sem conceito trinta e sete cursos.
Dos 113 cursos da região Sul, 67 eram de Universidades, para os quais o conceito modal foi 3, com 24 cursos.
Nesse tipo de organização, sete cursos ficaram sem conceito e os demais receberam conceito 2 ( 21 cursos ), conceito 4 ( 14 cursos ) e conceito 1 ( um curso ).
Os Centros Universitários tiveram cinco dos nove cursos participantes no conceito modal, o conceito 3.
Os demais quatro cursos receberam conceito 1 ( um curso ) e conceito 2 ( três cursos ).
Foram 37 os cursos vinculados a Faculdades na região Sul, e destes, 11 ficaram sem conceito.
Os restantes receberam conceito 2 ( 13 cursos, o conceito modal ), conceito 3 ( dez cursos ), conceito 1 ( dois cursos ) e conceito 4 ( apenas um deles ).
Na região Centro-Oeste, 25 dos 70 cursos eram de Universidades.
Destes, três cursos ficaram sem conceito e o conceito modal foi 4, com nove cursos.
Dos outros cursos, cinco foram avaliados com conceito 3, quatro com conceito 2, três com conceito 5 e um com o 1.
Os Centros Universitários contaram com nove cursos, três dos quais não receberam conceito.
Dos seis cursos restantes, distribuíram -se igualmente nos conceitos 2 e 3, conceitos modais.
Dos 36 cursos em Faculdades, 19 ficaram sem conceito e 12 receberam o conceito modal entre os com conceito, 2.
Dois cursos receberam conceito 1 e três o conceito 3.

% matched lemmas: 
\end{textsample}
